\documentclass[11pt]{article}
\usepackage[T1]{fontenc}
\usepackage{lmodern}
\usepackage[letterpaper,margin=1in]{geometry}
\usepackage{amsmath,amssymb,amsthm,mathtools,microtype}
\usepackage[hidelinks]{hyperref}
\hypersetup{pdftitle={Odd azimuthal graviton pairs in the resonant dynamics of global AdS4}}
\newtheorem{theorem}{Theorem}
\newtheorem{lemma}[theorem]{Lemma}
\newcommand{\R}{\mathbb R}
\newcommand{\Lie}{\mathcal L}
\newcommand{\eps}{\varepsilon}
\newcommand{\Sym}{\operatorname{Sym}}
\newcommand{\dd}{\,\mathrm d}
\numberwithin{equation}{section}
\title{Odd azimuthal graviton pairs in the resonant dynamics of global \(\mathrm{AdS}_4\)}
\author{}
\date{}
\begin{document}
\maketitle
\begin{abstract}
A nonlinear classical state invariant under rotation by \(\pi\) cannot
develop an odd azimuthal field. The corresponding quantum symmetry
allows two odd gravitons. We show that the constraint-eliminated
quartic resonant Hamiltonian of pure Einstein gravity in
conformal-Dirichlet global \(\mathrm{AdS}_4\) actually generates such a
pair. The AdS boost Ward identity must be applied after both the
Hamiltonian and the charges are transformed to normal form. Its
restriction to three two-graviton states, together with a known
nonzero classical Einstein secular residue, gives a quantitative
lower bound for transfer into the odd azimuthal pair \((3,5)\).
We construct a unitary Wick dynamics on positive Fock space
through finite energy-grade blocks and an exactly fixed-energy-grade,
fixed-axial-charge packet. At time \(s\sqrt h\) in the resonant
dynamics, this packet has an odd-pair probability
\(h s^2 |q_{35}|^2+o(h)\), \(q_{35}\ne0\), alongside stimulated even
daughters and an explicitly constrained Gaussian stock fluctuation.
The result is for the reduced quartic quantum theory and matched
mode readout, not an all-orders quantization of Einstein gravity.
\end{abstract}

\section{Setting and statements}\label{sec:intro}

Set the AdS radius to one. The conformal boundary is the marked round
cylinder \((\R_t\times S^2,-\dd t^2+\gamma_{S^2})\). We impose a smooth
center and conformal Dirichlet falloff, not a finite-radius wall. In a
positive canonical free mode frame write \(z_j\) for the complex
amplitude. The index \(j\) includes the polar or axial master type,
\(\ell\ge2\), \(-\ell\le m\le\ell\), and radial overtone \(p\ge0\).
The two free frequency towers are
\begin{equation}\label{eq:towers}
 \omega_{\rm pol}=\ell+1+2p,\qquad
 \omega_{\rm ax}=\ell+2+2p .
\end{equation}
Thus \(H_2=\sum_j\omega_j|z_j|^2\) and
\(J_z=\sum_jm_j|z_j|^2\), with
\(i\dot z_j=\partial_{\bar z_j}H\). We denote the unique polar mode
\((\ell,m,p)=(j,j,0)\) by \(C_j\) when its label is a single integer.
Its defect \(\omega-m\) equals one. \emph{Odd} below always refers to
the marked azimuthal number, not to the axial master type.

Let \(H_{\mathrm{res}}\) be the quartic stationary normal-form
coefficient of the genuine constrained Dirichlet Einstein
Hamiltonian (Section~\ref{sec:ward}). On bosonic Fock space let \(a_j\) be
the unit-CCR annihilators, \(N_j=a_j^\dagger a_j\), and
\[
 K=\sum_j\omega_jN_j,\qquad M=\sum_jm_jN_j .
\]
Write \(\widehat H_h\) for the normal ordering of \(H_{\rm res}\) with
\(z_j\mapsto\sqrt h\,a_j\), \(\bar z_j\mapsto\sqrt h\,a_j^\dagger\), and
\(U_h(\tau)=\exp(-i\tau\widehat H_h/h)\). Here \(h>0\) is the
dimensionless semiclassical parameter of the resonant model, and
\(\tau\) is its slow time. The unit-\(h\) matrix is
\(V=\widehat H_1|_{\text{two chiral particles}}\). In its joint
\(K=10,M=8\) subspace use the normalized states
\[
 A=|C_2 C_6\rangle,\qquad B=|C_3 C_5\rangle,\qquad
 C=|(C_4)^2\rangle=\frac{(a_4^\dagger)^2}{\sqrt2}\Omega .
\]
\begin{theorem}[An odd Einstein graviton channel]\label{thm:main}
The quartic resonant operator defines a self-adjoint direct sum on
the physical positive-mode Fock space, preserving \(K,M\); the chiral
sector generated by the \(C_j\) is invariant. In a canonical real
boost-ladder phase convention its two-particle block obeys
\begin{equation}\label{eq:key}
 V_{AC}=\frac{3\sqrt{14}}4 V_{AB}-\frac{\sqrt{10}}2 V_{BC},
 \qquad
 |V_{BA}|^2+|V_{BC}|^2\ \ge\ \frac{8}{83}|V_{AC}|^2>0.
\end{equation}
The absolute bound is invariant under rephasing the individual modes.
In particular the transition to \((3,5)\) from at least one of the
two even pairs \((2,6)\) and \((4,4)\) is nonzero.
\end{theorem}

For an even stock \(z=(z_2,z_4,z_6)\) with \(I_i=|z_i|^2>0\),
put \(X=-i\partial_{\bar z}H_{\rm res}\) and set
\begin{equation}\label{eq:parameters}
 D=\operatorname{diag}(I_2,I_4,I_6),\quad
 L=\begin{pmatrix}1&1&1\\1&2&3\end{pmatrix},\quad
 \mu_i=2\operatorname{Re}(\bar z_i X_i(z)),\quad
 q_{35}=V_{BA}z_2z_6+\frac{V_{BC}}{\sqrt2}z_4^2 .
\end{equation}
Select integers \(c_h\in\mathbb Z^2\) with \(c_h-LI/h=O(1)\).
Let \(\Pi_h\) project onto
\((K,M)=(c_{h,1}+2c_{h,2},2c_{h,2})\), and form the normalized
vector from \(\Pi_h D_{\rm W}(z/\sqrt h)\Omega\), where \(D_{\rm W}\)
is the usual coherent displacement.
\begin{theorem}[Sharp-charge coincidence and recoil]\label{thm:packet}
There exist stocks as above for which \(q_{35}\ne0\),
\(\mu\ne0\), and \(X_8(z)\ne0\). Evolve the projected vector
by \(U_h(s\sqrt h)\) with fixed \(s\ne0\). The probability that both
\(N_3\) and \(N_5\) are nonzero is
\begin{equation}\label{eq:rare}
 h s^2 |q_{35}|^2+o(h).
\end{equation}
The first even daughters \(C_8,C_{10}\) have limiting independent
Poisson counts, of means \(s^2|X_8(z)|^2\) and \(s^2|X_{10}(z)|^2\).
They are independent of the limiting stocked fluctuation:
\begin{equation}\label{eq:gauss}
 \frac{hN_{\rm stock}-I}{\sqrt h}
 \ \Longrightarrow\
 {\cal N}\big(s\mu,D-DL^{\mathsf T}(LDL^{\mathsf T})^{-1}LD\big).
\end{equation}
The covariance has rank one in the direction \((1,-2,1)\).
Conditioning on the odd event gives the same leading even and
stocked laws, with a shift of \((2,4)\) in the exact stock
\((K-M,M/2)\) budget before accounting for any other daughters.
The rare coefficient and limiting laws are uniform for each
physical mode cutoff containing a fixed finite set used by the proof.
\end{theorem}
The classical Einstein quartets are established by Dias,
Horowitz and Santos and by Dias and Santos \cite{dhs,ds}.
Algebraic selection, weakly nonlinear dynamics, finite quantum
resonance blocks and maximally rotating sectors have substantial
prior work \cite{evninniv,evninreview,ep,cde,hk}. Our contribution
is the transformed gravitational boost Ward identity as applied
to the concrete block, and the exact-charge quantum odd-pair
prediction forced by an actual gravitational secular residue.

\section{The constrained quartic and the boost Ward identity}\label{sec:ward}
\subsection{Coefficient definition and the classical Einstein input}

Use the Einstein--Hilbert action with cosmological term, the
Dirichlet boundary terms and the AdS normalization of its charges.
Proper diffeomorphisms vanish as boundary transformations and carry
no charge. After solving the constraints and quotienting those
diffeomorphisms, use the reduced canonical charge algebra at AdS.
The formulas below are for its Taylor coefficients: they do not
posit a global smooth canonical chart on nonlinear Cauchy data.
For finitely many external linear modes the intermediate-order
exchange is computed by the forced linear equation at each
signed input frequency and allowed angular label, with a smooth-center
solution and the required inhomogeneous Dirichlet metric lift. This
includes the exceptional low angular harmonics, and does not assign
independent homogeneous Robin data to sourced polar masters.
The metric prescription is that of \cite[Section II and Appendix A]{ds}.
Corner compatibility of an exact all-order metric family is a
different problem \cite{soutetre}. A coefficient computed with an
intermediate arbitrary spectral cutoff is not the definition.

Write \(H=H_2+H_3+H_4+\cdots\) in a parity- and rotation-equivariant
reduced canonical frame (the finite-order Darboux construction is
equivariant by averaging its primitive over this compact group).
Let \(\mathcal L=\{H_2,\cdot\}\) and denote its grade \(k\)
projection by \(\Pi_k\), with eigenvalue \(ik\) and free period \(2\pi\).
A signed degree-\(d\) monomial of grade \(k\) has parity
\((-1)^{k-d}\), since both master types in \eqref{eq:towers}
have parity \((-1)^{\omega-1}\). Thus no parity-even cubic
has grade zero. Choose the cubic generator \(F_3\) on nonzero
grades by \(\{H_2,F_3\}=-H_3\), and a quartic generator \(F_4\)
eliminating the nonzero grades at degree four. Set
\begin{equation}\label{eq:nf}
 H_{\rm res}=\Pi_0\big(H_4+\tfrac12\{H_3,F_3\}\big),
 \qquad H'=H_2+H_{\rm res}+O(5).
\end{equation}
The algebraic inverse of \(\mathcal L\) divides by nonzero detuning.
The physical second-order metric Green solve, in a different
space, evaluates the constrained exchange before taking the
third-order resonant projection. In particular that exchange,
including the low-multipole pieces, is not a bare quartic vertex.

The reduced one-period return \(P\) has linear part the
identity and no quadratic part, by integer frequencies and the
cubic parity rule. Its cubic jet is \(P_3=2\pi X\) in the
normal frame. It is invariant under an equal-cut tangent-identity
change of coordinates: conjugating
\(\mathrm{Id}+P_3+O(4)\) by such a map leaves its first
nontrivial jet. On any normalizable marked boundary observable
\(O\), \(O(P(z))-O(z)=\dd O_0 P_3(z)+O(4)\).
Thus the gauge-invariant secular response in a vacant mode
cannot be made or cancelled by a periodic change of variables.
Homogeneous second-order \emph{radiative} additions cannot
enter this coefficient either, since their coupling at
third order would be a stationary cubic vertex, which is absent.
A frequency adjustment in an occupied linear mode does not
change the distinct vacant mode. The forced low multipoles of
the Einstein solve are not additional free radiation.

In \cite[Eqs.~(124), (132), (134), arXiv v1]{ds}, the \(C_4,C_6\)
seed has a nonzero \((\omega,m)=(3,2)\) secular boundary-stress
residue of amplitude proportional to the square of \(C_4\) times
\(C_6\). These quantum numbers select a unique physical mode.
Of the cubic contractions from that seed only \(z_4^2\bar z_6\)
has those weights. The same computation has a nonzero \((9,8)\)
residue from \(z_6^2\bar z_4\). Because a nonzero linear mode
normalization does not change nonvanishing, this gives
\begin{equation}\label{eq:inputs}
 H_{\rm res}\supset
 \Lambda z_4^2\bar z_2\bar z_6+\bar\Lambda\bar z_4^2z_2z_6,
 \qquad \Lambda\ne0;\qquad
 [z_6^2\bar z_4]X_8\ne0 .
\end{equation}
The reciprocal \(C_6\) secular daughter of a \(C_2,C_4\) seed was
found independently in \cite[discussion around Eq.~(8)]{dhs}.
We use no numerical value for the canonically normalized
\(\Lambda\); the classical papers use metric amplitudes.

\subsection{Transforming the equal-cut charges}

Boundary conformal Killing fields with compensating radial
diffeomorphism preserve the fixed boundary representative.
They are physical Hamiltonian symmetries, not proper gauge.
The bracket and integrable surface charges are described in
\cite[Secs.~2.2--2.3]{him}, \cite[Secs.~4.1--4.3]{ps}, and
\cite[Sec.~3.1]{cm}. The boundary dimension is three and there
is no Weyl anomaly. Normalize all charges to zero on AdS.
Their Hamiltonian vector fields have a fixed point there, so
their Poisson bracket has no residual central constant.
A separately marked readout does not modify the action or its
space of boundary symmetries.

Take an equal-cut complex boost \(B=B_2+B_3+B_4+\cdots\)
with \(\{H,B\}=i\lambda B\), \(\lambda=\pm1\), after factoring
the explicit linear time dependence of the conserved charge.
Apply the same normal-form transformations to \(B\), not only
to \(H\). Then \(B'_2=B_2\), \(B'_3=B_3+\{B_2,F_3\}\), and
\[
 B'_4=B_4+\{B_3,F_3\}
       +\tfrac12\{\{B_2,F_3\},F_3\}+\{B_2,F_4\}.
\]
The degree-four bracket gives
\[
  (\mathcal L-i\lambda)B'_4=-\{H_{\rm res},B_2\}.
\]
There is no \(\{H'_3,B'_3\}\), since \(H'_3=0\).
By Jacobi the right side has grade \(\lambda\), while
\(\Pi_\lambda\) annihilates the left side. Therefore
\begin{equation}\label{eq:ward}
 \{H_{\rm res}, B_2\}=0 .
\end{equation}
No vanishing of \(B'_4\) itself is asserted. This is a
Ward identity for constrained quartic coefficients before
spectral compression, not for a nonlinear boost on a finite mode
cutoff or an all-order quantum charge. Rotation invariance also
gives \(\{H_{\rm res},J_z\}=0\).


\section{The physical ladder and a nonzero odd matrix element}\label{sec:matrix}

On the boundary cylinder let \(Y=\sin\theta e^{i\phi}\). The
conformal Killing field
\[
 \xi_+=e^{-it}\big(Y\partial_t+i\nabla^A Y\partial_A\big)
\]
obeys \(\Lie_{\xi_+}g_{(0)}=-2i e^{-it}Y g_{(0)}\).
In a Fefferman--Graham defining coordinate \(x\) this variation is
compensated by the normal component
\(\xi_+^x=-i e^{-it}Yx+O(x^2)\) and the associated subleading
tangential extension. Put \(\xi_-=\bar\xi_+\). The identities
\(\nabla_A\nabla_B Y=-Y\gamma_{AB}\),
\(|Y|^2+\nabla Y\cdot\nabla\bar Y=2\) and
\(\bar Y\nabla Y-Y\nabla\bar Y=2i\partial_\phi\)
give \([\xi_-,\xi_+]=4i(-\partial_t+\partial_\phi)\);
\(\xi_+\) raises both \(\omega\) and \(m\) by one on
\(e^{-i\omega t+im\phi}\). The positive one-particle sector
is invariant under the infinitesimal action, since there
are no zero or adjacent opposite-frequency states.
The real Killing actions are skew-symmetric in the positive
free inner product; on the finite-mode core use
\(K_+=\tfrac12\Lie_{\xi_+}\) and
\(K_-=-\tfrac12\Lie_{\xi_-}=K_+^\dagger\).
They obey \([K_-,K_+]=2K_0\), where
\(K_0=(\omega+m)/2\) and \([K_0,K_\pm]=\pm K_\pm\).

\subsection{Closure without an assumed particle-number selection rule}

The positive defect \(\delta=\omega-m\) has lower bound one
by \eqref{eq:towers}. The defect-one chain consists exactly
of the \(C_j\), with lowest weight \(K_0=5/2\).
Since \(K_+\) and \(K_-\) preserve defect, an intertwining
argument gives the necessary absence of a \(3\)-to-\(1\)
resonance. The possible one-particle output has defect three.
That sector consists of the polar \(p=0,m=\ell-2\), axial
\(p=0,m=\ell-1\), and polar \(p=1,m=\ell\) chains of
mode labels. At fixed \(m=0,1,2,3,\ldots\) its dimensions are
\(1,2,3,3,\ldots\), with \(K_0=m+3/2\).
Since \(K_0>0\), the raising map is injective
(\(\|K_+u\|^2=\|K_-u\|^2+2\langle u,K_0u\rangle\)).
The adjoint relation implies that new lowest vectors arise
only at the three dimension jumps. The resulting lowest
weights \(3/2,5/2,7/2\) span the whole sector under raising.

Equation~\eqref{eq:ward} makes a \(3\)-to-\(1\) coefficient an
intertwiner from the symmetric cube of defect one to this
defect-three sector. All source weights are at least \(15/2\).
A nonzero image would contain a lowest-weight vector
by taking its lowest nonzero weight and lowering; the target
has none that high. Thus the coefficient vanishes. Its adjoint
cannot send a defect-one particle to three positive-defect
particles. A \(2\)-to-\(2\) term from defect-one inputs has
total defect two and keeps both outputs in defect one. Terms
with four creation or annihilation operators cannot conserve
positive free energy. These statements prove invariance and
number conservation on the chiral Fock space; no assumption
about every other sector is necessary. The normal-order map
intertwines the additive linear one-particle symmetry in
each coefficient: the quadratic symmetry generator is
number preserving, so commuting it through Wick monomials
acts on their separate inputs and outputs, without contractions.

\subsection{The three-channel matrix}

In a Bargmann realization of the lowest-weight chain,
\[
 K_-=\partial_w,\quad K_+=w^2\partial_w+5w,\quad
 \|w^n\|^2=\frac{n!}{(5)_n}.
\]
The index of \(C_j\) in this chain is \(j-2\).
On two symmetric particles, total degree four has the
basis of polynomials \(w_1^4+w_2^4\),
\(w_1^3w_2+w_1w_2^3\), \(w_1^2w_2^2\).
Their norms are \(n_A=1/\sqrt{35}\),
\(n_B=\sqrt{2/175}\), \(n_C=1/15\), respectively.
The lowest vectors \(1,(w_1-w_2)^2,(w_1-w_2)^4\)
under \(K_-^{(1)}+K_-^{(2)}\), raised to degree four,
give three distinct Casimir lines.
In the displayed unnormalized basis their coordinates, up
to a common scalar in each row, are
\begin{equation}\label{eq:lines}
 (28,70,90),\qquad (4,-1,-6),\qquad (1,-4,6).
\end{equation}
For example \(K_+^{(1)}+K_+^{(2)}\) acting once on the
second lowest vector gives \(7(w_1^3+w_2^3-w_1^2w_2-w_1w_2^2)\).
The distinct Casimir eigenvalues and the three independent
vectors exhaust this block. The Hermitian operator \(V\)
commutes with the Casimir, hence is a sum of real multiples
of their orthogonal projections. For each row \((a,b,c)\) in
\eqref{eq:lines}, \(ac=(9/2)ab-bc\).
Multiplying by the basis norms and normalizing each line
gives the entrywise first identity in \eqref{eq:key}.
Cauchy's inequality gives the factor \(8/83\), since
\((3\sqrt{14}/4)^2+(\sqrt{10}/2)^2=83/8\).
The source quartet \eqref{eq:inputs} gives
\(V_{AC}=\sqrt2\,\Lambda\ne0\), using the normalized
identical-particle vector \(C\). This proves the nonzero
odd channel. The argument would not go through by merely
listing frequency-allowed terms.

Classically the fixed-point subspace of rotation by \(\pi\)
is invariant even in the constrained normal form. In it,
every odd azimuthal \emph{amplitude} vanishes. A quantum state
in the even eigenspace of this rotation only has even
\emph{total} azimuthal charge; in particular the pair \(B\)
is allowed. Both \(A\) and \(C\) have the same conserved
charges as \(B\). This is the difference tested here.


\paragraph{Finite-time block witness.}
There is also an exact non-asymptotic test in the three-channel
matrix alone. Prepare the equally weighted incoherent mixture of
\(A\) and \(C\) in the unit-\(h\) block. By unitarity its
\(B\)-probability under \(e^{-itV}\) is
\((1-|\langle B,e^{-itV}B\rangle|^2)/2\).
The weights of \(B\) on the three normalized lines in
\eqref{eq:lines} are \(70/143,1/55,32/65\).
If \(E_\ell\) are the eigenvalues on these lines, the time mean
on \([0,T]\) is
\[
 \sum_{\ell<j} p_\ell p_j
     \big[1-\operatorname{sinc}(T(E_\ell-E_j))\big],
 \qquad \operatorname{sinc}(u)=\sin(u)/u .
\]
The spectral diameter is at least \(2|V_{AC}|\): subtract the
spectral midpoint and use \(|V_{AC}|\le
\|V-\tfrac12(E_{\max}+E_{\min}){\bf 1}\|\).
At \(T\ge |V_{AC}|^{-1}\), the extremal pair contributes at
least \(7/1573\) even if two eigenvalues are equal.
In the infinite-time mean, distinct eigenvalue groups
give at least \((1/55)(54/55)=54/3025\).
These are bounds for the retained block dynamics, not for
full Einstein evolution at long times.

\section{Exact dynamics of the quartic Wick model}\label{sec:fock}

For each integer \(k\), \(\ker(K-k)\) is finite dimensional:
the free spectrum is bounded below by three, has finite
multiplicities, and no mode of frequency greater than \(k\)
can be occupied. Normal ordering gives a Hermitian matrix
in every such block and preserves \(M\) termwise. Define
\(\widehat H_h\) as their self-adjoint direct sum with domain
\(\sum_k\|\widehat H_{h,k}\psi_k\|^2<\infty\).
Finite block sums are a core. This is positive Hilbert-space
unitarity, not a claim that the quartic energy is bounded below.
A finite physical mode cutoff is a Hermitian compression of
these same coefficients. Once the cutoff contains all modes
needed for grade \(k\), it agrees exactly on that block.
The strong convergence of \(U_h\), at fixed \(h\), therefore
follows by first retaining finitely many blocks and using the
unitary norm bound on the tail. This construction involves no
vacuum-energy subtraction from an infinite non-Wick ordering.

For later uniformity consider a coherent displacement on a
finite set of modes of maximum frequency \(W\), followed by a
polynomial in finitely many creation operators of frequencies
\(\le W\) acting on the displaced vacuum in the corresponding
Fock frame. Wick-expand after conjugating by the displacement.
In any nonzero term every annihilated leg either contracts
with that finite stock or with an explicitly present finite
excitation. Conservation of the positive grade bounds every
new created frequency by \(3W\) (in the chiral \(2\)-to-\(2\)
case \(2W\) already suffices). Thus any fixed finite sequence
of applications reaches only finitely many labels, and the
bounds on these particular vectors hold identically for all
larger cutoffs. It follows in particular that the trial vectors
used below lie in the full graph domain: their Wick images
are finite sums of displaced finite excitation vectors, whose
norms are finite; checking against every finite \(K\) sector
identifies the direct-sum operator. No uniform norm estimate
on an entire growing block is needed.

\subsection{Control with a single pump}

An independent sharp-number control removes all phase
preparation. Take exactly \(n_h\) quanta of \(C_4\) with
\(h n_h\to I_4>0\) and all others empty. By chiral invariance
and the two-to-two rule the classical restricted amplitude
returns to \(C_4\). A quantum departure can go to \(A\) or
\(B\), both with the pump occupation lowered by two.
Subtract the scalar diagonal value on the seed from the
generator \(\widehat H_h/h\), calling the result \(G_h\).
In the first two fixed reachability steps from the seed,
the pure-pump diagonal difference is \(O(1)\), since it
is \(h\) times the difference of \(n(n-1)\) for
\(n=n_h+O(1)\). Every off-diagonal term has at most two
pump legs: three pump momenta with one distinct leg are
excluded by the two-to-two marked selection, and the
one-to-three terms on this domain have been ruled out.
Thus such matrix elements are \(O(1)\), with the two pump
ladder factors at most \(O(h^{-1/2})\).
The finite second neighborhood yields
\(\|G_h^2|n_h;C_4\rangle\|=O(1)\).
Taylor's spectral bound gives an error \(O(\tau^2)\)
after subtracting the first derivative of the vector.
The exact daughter-pair amplitude into \(C_2,C_6\) with
\(n_h-2\) remaining \(C_4\) quanta is consequently
\[
 -i\tau\Lambda h\sqrt{n_h(n_h-1)}+O(\tau^2).
\]
Its probability at \(\tau=s\sqrt h\) is
\(h s^2|\Lambda|^2I_4^2+o(h)\).
This separately verifies quantum departure from a
single-mode classical invariant ray. It does not assert
that this pump alone has necessarily \(V_{BC}\ne0\).


\section{The sharply conditioned packet}\label{sec:packet}

\subsection{Choice of stock and the classical test}

Write \(z_2,z_6>0\) and \(a=z_2z_6>0\), and choose
\(z_4^2=\pm i\sqrt2\,a\).
All entries of \(V\) are real in the ladder phases. Thus
\[
 |q_{35}|^2=a^2\big(|V_{BA}|^2+|V_{BC}|^2\big)
       \ge (8/83)a^2|V_{AC}|^2 .
\]
Up to diagonal action terms, the only resonance within the
three stocked modes is \eqref{eq:inputs}. It gives the
action derivative
\[
 \mu=2\operatorname{Im}(\Lambda z_4^2\bar z_2\bar z_6)
                  (1,-2,1),
\]
nonzero for either phase, with opposite signs. A free
one-creation variation from the even stock can only use
\(j_{\rm out}=j_1+j_2-j_3\) in the chiral sector; hence the
only empty outputs in \(X(z)\) are \(C_8,C_{10}\).
The two terms at label \(8\) have forms
\(\Gamma z_6^2\bar z_4+\Gamma' z_4z_6\bar z_2\).
Their squared magnitude averaged over the two displayed
phase choices is the sum of the squared term magnitudes,
and \(\Gamma\ne0\) by \eqref{eq:inputs}. At least one choice
has \(X_8\ne0\), proving all three strict conditions in
Theorem~\ref{thm:packet}. The positive-defect mode closure
also implies \(X_3=X_5=0\).

\subsection{Two-step expansion in rescaled time}

Let \(\tau=s\sqrt h\), and before charge projection write
the evolved coherent state in its fixed frame
\(D_{\rm W}(z/\sqrt h)\). Remove the scalar phase
from the constant term of the shifted generator.
Normal ordering gives the exact rescaled generator
\begin{equation}\label{eq:gen}
 \mathcal L_z+\sqrt h\,\mathcal Q_z
       +h\,\mathcal C_z+h^{3/2}\mathcal V .
\end{equation}
Here \(\mathcal L_z,\mathcal Q_z,\mathcal C_z,\mathcal V\)
are the linear, quadratic, cubic and quartic normal-ordered
Taylor terms at \(z\), written in unit-CCR fluctuation
operators. In particular
\(\mathcal L_z=i\sum_j(X_j a_j^\dagger-\bar X_j a_j)\)
and its evolution is the coherent displacement \(D_{\rm W}(sX)\).
In the further frame of that displacement the first correction
to the vacuum is generated by
\(D_{\rm W}(-sX)\mathcal Q_zD_{\rm W}(sX)\).
Translation does not change the two-creation coefficients
of this quadratic term. Its component with both modes \(3,5\)
is exactly
\begin{equation}\label{eq:coefficient}
 -i s\,q_{35}\,a_3^\dagger a_5^\dagger\Omega .
\end{equation}
In fact the \(B\leftarrow A\) term multiplies \(z_2z_6\), and the
\(B\leftarrow C\) term multiplies \(z_4^2/\sqrt2\)
by identical-particle normalization. In this frame the
displacement is entirely even.

To obtain a small enough error after sharp conditioning, let
\(\chi_0(s)=D_{\rm W}(sX)\Omega\) and recursively set
\[
 i\dot\chi_j=\mathcal L_z\chi_j+
     \mathcal Q_z\chi_{j-1}+\mathcal C_z\chi_{j-2}
         +\mathcal V\chi_{j-3}, \qquad j=1,2 ,
\]
with initial value zero and negative-index terms zero.
After the \(D_{\rm W}(-sX)\) displacement, these functions
are finite polynomials of degree bounded independently of \(h\),
on finite frequency sets by Section~\ref{sec:fock}.
Their coefficients stay bounded on compact \(s\)-intervals,
since translating the fixed Wick polynomials produces
finite polynomial functions of \(s\). Substitution of
\(\chi_0+\sqrt h\,\chi_1+h\chi_2\) into \eqref{eq:gen}
leaves a norm residual \(O(h^{3/2})\). Duhamel's formula
then gives the same norm bound for the vector error.
This calculation is valid against the direct-sum
unitary on the stated graph domain, and has identical
estimates for every sufficiently enlarged physical cutoff.
Denote the combined leading displacement by
\(D_{\rm W}(y/\sqrt h)\) with \(y=z+s\sqrt h\,X(z)\);
the phase from composing the two displacements is
common to the comparison vectors and is irrelevant.

\subsection{A charge-projector bound for finite excitations}

Put \(D_h=D_{\rm W}(y/\sqrt h)\) and denote a single joint
\(K,M\) projector by \(\Pi\). For a bounded finite linear
combination \(\eta\) of finite-excitation basis states on
fixed finite mode sets there is a uniform charge bound
\begin{equation}\label{eq:projector}
 \|\Pi D_h\eta\|\le C_\eta\sqrt h .
\end{equation}
To see this, integrate the diagonal matrix element of
\(\exp(i\theta K+i\phi M)\) over the physical charge torus.
The stocked coherent part of its absolute value is
\[
 \exp\!\left[-\frac{1}{h}
  \sum_{j=2,4,6}|y_j|^2(1-\cos(\omega_j\theta+j\phi))\right].
\]
Finite excitations, via the Weyl formula and normal ordering,
multiply this by a bounded-degree polynomial in
\((e^{i(\omega_j\theta+j\phi)}-1)y_j/\sqrt h\) and conjugates
with bounded coefficients. The other coherent displacements
are bounded and the unexcited spectator factors have norm
at most one. A uniform fraction of the displayed exponential
absorbs the polynomial. The three stock rows are the image
of the primitive columns \((1,1),(1,2),(1,3)\) under the
lattice map to \(K,M\); their weights span rank two with
index two on the physical torus. There are therefore just
two common saddles, locally Gaussian in two independent
directions, and away from them the modulus is exponentially
small. Their integrals are \(O(h)\) for the squared norm,
uniformly in the projector target. This proves
\eqref{eq:projector}. In particular the \(h\chi_2\) correction
after sharp projection and normalization will cost \(O(h)\),
not \(O(\sqrt h)\).

\subsection{Primitive lattice inversion and the rare event}

In the primitive coordinates \((K-M,M/2)\), the stocked
coherent occupations have independent Poisson laws with means
\(|y_j|^2/h\) and charge columns \(L_j\) in
\eqref{eq:parameters}. The even daughters' columns are
also integral: \(C_j\) with even \(j\) carries \((1,j/2)\).
Since \(L\mu=0\), the stocked charge mean changes only
by \(O(1)\) when replacing \(z\) by \(y\), at this scale.
Fourier inversion on the primitive torus for a stocked charge
target \(b_h=LI/h+O(1)\) reads
\[
 \frac1{(2\pi)^2}\int_{[-\pi,\pi]^2}
 \exp\!\left(\sum_j\frac{|y_j|^2}{h}
             (e^{i L_j\cdot\theta}-1)-i b_h\cdot\theta\right)
 \dd^2\theta
 =\frac{h}{2\pi\sqrt{\det(LDL^{\mathsf T})}}\,(1+o(1)).
\]
Indeed only \(\theta=0\) on this torus is a common saddle;
with \(\theta=\sqrt h\,\eta\) its positive quadratic
matrix is \(LDL^{\mathsf T}\). This estimate is also
uniformly \(O(h)\) for \emph{all} targets by dropping
the unit-modulus target factor and bounding the same
Gaussian modulus. By convolution with the bounded-mean
even daughter Poisson counts and domination, the
leading coefficient is unchanged when any fixed daughter
charges, or the fixed shift \((2,4)\), are included.
This also shows that the unshifted initial projector has
norm asymptotic to a positive constant times \(\sqrt h\).

Let \(E_B\) project onto simultaneous \(N_3>0,N_5>0\).
It commutes with \(\Pi_h\). The initially unprojected state
has no odd displacement, and the corresponding \(E_B\) component
of the first correction is precisely \eqref{eq:coefficient}
in the frame of \(D_h\); all other first corrections have
degree at most two and cannot have both modes.
The first real two-creation vector has unit norm before
displacement and projection. Its two particles carry exactly
\((K-M,M/2)=(2,4)\). Conservation allows projecting after
evolving the initial unprojected coherent state and dividing
by the initial projected norm. The residual \(O(h^{3/2})\)
and the projected \(h\chi_2\) term combined give \(O(h)\)
in the normalized vector by \eqref{eq:projector}. Therefore
the \(E_B\) vector has squared norm given to leading order by
\[
 h s^2 |q_{35}|^2\,
 \frac{\Pr_y[\text{even coherent charges }=c_h-(2,4)]}
      {\Pr_z[\text{initial coherent charges }=c_h]} .
\]
The ratio tends to one by the lattice calculation.
The interference with the \(O(h)\) amplitude error is
\(O(h^{3/2})\), proving \eqref{eq:rare}.

Finally insert the bounded characteristic factor
\(\exp(i u\cdot(hN_{\rm stock}-I)/\sqrt h)\) in the same
Fourier inversion, and fix any finite daughter-count vector.
After rescaling \(\theta\), the stocked integrand has limiting
Gaussian factor
\[
 \exp\!\left(-\tfrac12(u+L^{\mathsf T}\eta)^{\mathsf T}
                  D(u+L^{\mathsf T}\eta)+is\,u\cdot\mu\right).
\]
The absence of an extra phase in \(\eta\) uses \(L\mu=0\).
The shifted cosine bound gives an integrable Gaussian
majorant for fixed \(u\). Completing the square and dividing
by the unshifted local limit gives \eqref{eq:gauss}.
The daughter Poisson law does not change to leading order,
because the fixed daughter integers only shift a stocked
target by \(O(1)\); its unbounded tails are dominated by the
uniform \(O(h)\) stock bound. The full projected state
differs from the normalized leading coherent vector by
\(O(\sqrt h)\) in norm for bounded tests. On the odd event,
after its own normalization the same comparison holds with
the shifted target, since its leading amplitude is
of order \(\sqrt h\) and has an \(O(h)\) error. The
covariance in \eqref{eq:gauss} is the restriction of the
diagonal Gaussian precision to \(\ker L\), which is
spanned by \((1,-2,1)\). These independent leading laws
do not remove exact recoil: before taking the limit every
sample of the full state obeys its precise \(K,M\) budgets.
This completes Theorem~\ref{thm:packet}.

\section{Observable convention and scope}

Mode labels refer to the boundary-marked, reduced canonical
variables with the same normal-form reconstruction on both
stroboscopic cuts. The amplitude can be calibrated on the
linear physical perturbations by the free spectral projection,
rather than by occupation counts of a raw local nonlinear metric.
For example an integer-frequency filter is obtained from a smooth
compactly supported time aperture \(g\) with
\(\sum_{n\in\mathbb Z}g(t+2\pi n)=1\); multiply by the
chosen Fourier exponential. Together with the angular mark
and a nonzero calibrated radial response this separates linear
overtones. This is not a claim about higher-order
nonlinear detector errors at the same order as the rare
probability. A canonical change of the cubic-free chart leaves
the stationary quartic unchanged through this order:
a resonant quartic generator commutes with \(H_2\), and
higher-degree changes cannot contribute. Rephasing the
marked linear modes changes conventions, not the absolute
matrix bound or the transformed stock experiment. Arbitrary
nonlinear bare detectors need not share these calibrated counts.

The proof yields spontaneous odd azimuthal production
from an even stock in an exact positive quartic Wick model.
It does not supply every Einstein quartet numerically, a full
interacting quantum theory, an all-order AdS approximation,
a continuum UV construction of quantum gravity, or a
prediction of nonlinear collapse.

\begin{thebibliography}{99}
\bibitem{dhs}
O. J. C. Dias, G. T. Horowitz and J. E. Santos,
Gravitational turbulent instability of anti-de Sitter space,
\emph{Class. Quantum Grav.} \textbf{29}, 194002 (2012),
\href{https://arxiv.org/abs/1109.1825}{arXiv:1109.1825}.
\bibitem{ds}
O. J. C. Dias and J. E. Santos, AdS nonlinear instability:
breaking spherical and axial symmetries,
\emph{Class. Quantum Grav.} \textbf{35}, 185006 (2018),
\href{https://arxiv.org/html/1705.03065v1}{arXiv:1705.03065v1}.
\bibitem{him}
S. Hollands, A. Ishibashi and D. Marolf,
Comparison between various notions of conserved charges in
asymptotically AdS spacetimes,
\emph{Class. Quantum Grav.} \textbf{22}, 2881 (2005),
\href{https://arxiv.org/abs/hep-th/0503045}{arXiv:hep-th/0503045}.
\bibitem{ps}
I. Papadimitriou and K. Skenderis,
Thermodynamics of asymptotically locally AdS spacetimes,
\emph{JHEP} \textbf{08}, 004 (2005),
\href{https://arxiv.org/abs/hep-th/0505190}{arXiv:hep-th/0505190}.
\bibitem{cm}
G. Comp\`ere and D. Marolf,
Setting the boundary free in AdS/CFT,
\emph{Class. Quantum Grav.} \textbf{25}, 195014 (2008),
\href{https://arxiv.org/abs/0805.1902}{arXiv:0805.1902}.
\bibitem{soutetre}
L. Sou\^etre, Geometric reflective boundary conditions for
asymptotically Anti-de Sitter spaces (2025), \href{https://arxiv.org/abs/2507.20661}{arXiv:2507.20661}.
\bibitem{evninniv}
O. Evnin and R. Nivesvivat,
AdS perturbations, isometries, selection rules and the Higgs oscillator,
\emph{JHEP} \textbf{01}, 151 (2016),
\href{https://arxiv.org/abs/1512.00349}{arXiv:1512.00349}.
\bibitem{evninreview}
O. Evnin, Resonant Hamiltonian systems and weakly nonlinear dynamics
in AdS spacetimes,
\emph{Class. Quantum Grav.} \textbf{38}, 203001 (2021),
\href{https://arxiv.org/abs/2104.09797}{arXiv:2104.09797}.
\bibitem{ep}
O. Evnin and W. Piensuk, Quantum resonant systems, integrable and chaotic,
\emph{J. Phys. A} \textbf{52}, 025102 (2019),
\href{https://arxiv.org/abs/1808.09173}{arXiv:1808.09173}.
\bibitem{cde}
B. Craps, M. De Clerck and O. Evnin,
Time-periodicities in holographic CFTs,
\emph{JHEP} \textbf{09}, 030 (2021),
\href{https://arxiv.org/abs/2103.12798}{arXiv:2103.12798}.
\bibitem{hk}
D. Hunik-Kostyra and A. Rostworowski,
AdS instability: resonant system for gravitational perturbations
of \(\mathrm{AdS}_5\) in the cohomogeneity-two biaxial Bianchi IX ansatz,
\emph{JHEP} 	extbf{06}, 002 (2020),
\href{https://arxiv.org/abs/2002.08393}{arXiv:2002.08393}.
\end{thebibliography}
\end{document}
